\documentclass[11pt]{article}
\usepackage[utf8]{inputenc}
\usepackage[T1]{fontenc}
\usepackage{amsmath,amssymb}
\usepackage{booktabs}
\usepackage{graphicx}
\usepackage{hyperref}
\usepackage{xcolor}
\usepackage[margin=1in]{geometry}
\usepackage{natbib}
\usepackage{multirow}
\usepackage{array}

\title{Limits of Cross-Axis Redundancy Signals:\\A Multi-Model Analysis of KV Similarity\\and MLP Sparsity Tolerance}

\author{Anonymous}

\date{}

\begin{document}
\maketitle

\begin{abstract}
We investigate whether a single KV-cache redundancy signal can jointly drive KV eviction and MLP sparsity allocation---a hypothesis we term \emph{cross-axis coordination}. Concretely, we test whether EMA-based KV cosine similarity predicts per-token MLP sparsity tolerance across five 7--8B parameter models spanning two attention architectures (GQA, MLA proxy) and two training paradigms (instruction tuning, reasoning distillation). Our multi-model analysis yields a decisive negative result: of five models tested, only Llama-3.1-8B-Instruct exhibits a meaningful cross-axis signal ($r{=}0.47$, $\Delta_\text{pert}{=}{+}6.6\%$). Mistral-7B-Instruct ($r{=}0.13$), Qwen2-7B-Instruct ($r{=}0.22$), and two reasoning-distilled models ($r{\leq}0.19$) all fail to replicate the signal despite sharing GQA attention. Per-layer decomposition reveals the signal concentrates in layers 12--18 even on Llama (9/32 layers with $r{>}0.4$), and the downstream two-tier sparsity benefit is negligible ($\Delta_\text{acc}{=}{+}0.01$). We conclude that cross-axis coordination is a model-specific artifact rather than an architectural property. As a standalone contribution, we present CARC-KV, an $O(1)$ KV eviction method that achieves 12.8$\times$ speedup over $O(n)$ baselines at 32K sequence length with ${\leq}2\%$ accuracy loss.
\end{abstract}

\section{Introduction}

Efficient inference for large language models (LLMs) requires reducing computation along multiple axes: KV-cache memory, MLP activation cost, and attention overhead. Recent work has pursued each axis independently---H2O \citep{zhang2023h2o} and SnapKV \citep{li2024snapkv} for KV eviction, Deja Vu \citep{liu2023dejavu} for MLP sparsity, and StreamingLLM \citep{xiao2024streamingllm} for attention windowing. A natural question arises: \emph{can a single redundancy signal coordinate efficiency decisions across multiple axes simultaneously?}

The \emph{cross-axis redundancy hypothesis} proposes that tokens with high KV-cache similarity to a running EMA anchor carry low information content, and therefore tolerate higher MLP sparsity. If true, this would enable $O(1)$ joint KV eviction and MLP sparsity allocation from a single cosine-similarity computation---an appealing primitive for on-device inference.

Initial results on Llama-3.1-8B-Instruct appeared to confirm this hypothesis: Pearson $r{=}0.47$ between EMA KV similarity and MLP sparsity tolerance, with perturbation-based validation showing a $+6.6\%$ accuracy advantage for guided sparsity. However, scientific claims require replication across models and conditions.

In this paper, we report the results of a systematic replication study across five 7--8B models. Our findings are negative:

\begin{enumerate}
    \item The cross-axis signal does not replicate on Mistral-7B-Instruct ($r{=}0.13$), despite identical GQA head configuration (8 KV heads, group ratio 4).
    \item Qwen2-7B-Instruct, another GQA model with instruction tuning, shows weak signal ($r{=}0.22$) with high per-problem variance.
    \item Reasoning-distilled models (DeepSeek-R1-Distill-Llama-8B, $r{=}0.19$; DeepSeek-R1-Distill-Qwen-7B, $r{=}0.15$) fail regardless of base architecture.
    \item Even on Llama, the signal concentrates in middle layers (12--18) and yields negligible downstream benefit ($\Delta_\text{acc}{=}{+}0.01$).
\end{enumerate}

We contribute: (1) a rigorous multi-model negative result with per-layer and per-head decomposition, establishing that cross-axis coordination is model-specific rather than architecture-generic; (2) a hypothesis-testing protocol for evaluating cross-axis efficiency signals; and (3) CARC-KV, an $O(1)$ KV eviction method competitive with $O(n)$ baselines.

\section{Background and Hypothesis}

\subsection{KV-Cache Redundancy}

In autoregressive generation, each token's key and value vectors are cached for future attention computation. Long sequences---particularly reasoning traces with meta-cognitive tokens (``wait'', ``let me reconsider'')---produce substantial cache redundancy. KV eviction methods exploit this by removing low-importance entries.

Existing methods differ in complexity: H2O maintains cumulative attention scores ($O(n)$ per step), SnapKV uses observation-window scoring ($O(n)$), while StreamingLLM uses a fixed sliding window plus sink tokens ($O(1)$ but accuracy-limited).

\subsection{MLP Activation Sparsity}

Transformer MLP layers exhibit natural activation sparsity: for any given token, a subset of neurons produces near-zero activations. Contextual sparsity methods \citep{liu2023dejavu} exploit this by predicting which neurons to skip, but require per-token sparsity predictors.

\subsection{Cross-Axis Coordination Hypothesis}

We define the cross-axis coordination hypothesis as follows:

\begin{quote}
\emph{If a token's KV representation is highly similar to a running EMA anchor (high redundancy), then that token carries low information content and tolerates higher MLP sparsity.}
\end{quote}

The causal chain is: high KV similarity $\rightarrow$ low surprisal $\rightarrow$ high MLP sparsity tolerance. This predicts a positive correlation between EMA KV cosine similarity and the maximum MLP sparsity a token can tolerate without accuracy degradation.

\subsection{EMA Anchor Similarity}

We define the EMA anchor $\mathbf{a}_t$ at position $t$ as:
\begin{equation}
\mathbf{a}_t = \alpha \cdot \mathbf{k}_t + (1 - \alpha) \cdot \mathbf{a}_{t-1}, \quad \mathbf{a}_0 = \mathbf{k}_0
\end{equation}
where $\mathbf{k}_t$ is the key vector and $\alpha = 0.1$. Token redundancy is:
\begin{equation}
s_t = \cos(\mathbf{k}_t, \mathbf{a}_{t-1})
\end{equation}
This is $O(1)$ per token---a single dot product and normalization, independent of sequence length.

\section{Experimental Protocol}

\subsection{Models}

We evaluate five models spanning two attention architectures and two training paradigms (Table~\ref{tab:models}).

\begin{table}[h]
\centering
\caption{Models evaluated. All are 7--8B parameters with GQA attention. The cross-axis signal is measured by Pearson $r$ between EMA KV similarity and MLP sparsity tolerance, and by perturbation accuracy differential $\Delta_\text{pert}$ at 55\% sparsity.}
\label{tab:models}
\small
\begin{tabular}{@{}lcccccc@{}}
\toprule
Model & KV Heads & Attn & Training & $r$ & $\Delta_\text{pert}$@55\% & Verdict \\
\midrule
Llama-3.1-8B-Instruct & 8 & Full & Instruction & 0.471 & +6.6\% & PASS \\
Mistral-7B-Instruct-v0.3 & 8 & SWA & Instruction & 0.13 & +0.8\% & FAIL \\
Qwen2-7B-Instruct & 4 & Full & Instruction & 0.22 & +1.2\% & FAIL \\
DeepSeek-R1-Distill-Llama-8B & 8 & Full & Reasoning & 0.19 & $-$0.8\% & FAIL \\
DeepSeek-R1-Distill-Qwen-7B & 4 & Full & Reasoning & 0.147 & $-$2.2\% & FAIL \\
\bottomrule
\end{tabular}
\end{table}

\subsection{Evaluation Tasks}

All experiments use 20 mathematical reasoning problems. For each problem, we collect per-token KV representations, compute EMA similarity scores, and measure MLP sparsity tolerance via perturbation analysis.

\subsection{Metrics}

\paragraph{Pearson $r$.} Correlation between EMA KV cosine similarity $s_t$ and per-token MLP sparsity tolerance (maximum sparsity before output deviation exceeds threshold). We require $r > 0.35$ for a PASS verdict.

\paragraph{Perturbation differential $\Delta_\text{pert}$.} At a fixed sparsity level (55\%), we compare accuracy when applying higher sparsity to high-similarity tokens versus low-similarity tokens. Positive $\Delta_\text{pert}$ confirms the hypothesis direction.

\paragraph{Per-problem consistency.} We report the range of per-problem $r$ values; high variance indicates an unreliable signal.

\paragraph{KS statistic.} Kolmogorov--Smirnov test on MLP activation distributions between high- and low-similarity token groups at various sparsity levels (L0, L2, L4).

\section{Results}

\subsection{Cross-Model Replication Failure}

Table~\ref{tab:models} summarizes the central result: only Llama-3.1-8B-Instruct exhibits a cross-axis signal meeting our threshold ($r{=}0.471$, $\Delta_\text{pert}{=}{+}6.6\%$). The remaining four models all fail.

\paragraph{Mistral-7B-Instruct.} Despite sharing identical GQA head configuration with Llama (8 KV heads, group ratio 4), Mistral yields $r{=}0.13$ with per-problem range $[-0.05, 0.28]$---some problems show \emph{negative} correlation. $\Delta_\text{pert}{=}{+}0.8\%$ is within noise. The key difference is Mistral's sliding window attention (SWA, window size 4096), which limits the attention span and causes the EMA anchor to lose representativeness beyond the window boundary. This demonstrates that even identical head configurations are insufficient---the attention mechanism itself matters.

\paragraph{Qwen2-7B-Instruct.} With 4 KV heads (group ratio 7) and full attention, Qwen2 shows $r{=}0.22$ and $\Delta_\text{pert}{=}{+}1.2\%$. The higher group ratio forces each KV head to serve more query heads, producing more generic KV representations that lose token-specific information. Per-problem $r$ ranges from 0.02 to 0.38, indicating an unstable signal.

\subsection{Training Paradigm Confound}

The confound control experiments reveal that training paradigm, not architecture, is the dominant factor.

\begin{table}[h]
\centering
\caption{Confound control: same GQA architecture, different training paradigms.}
\label{tab:confound}
\small
\begin{tabular}{@{}llccc@{}}
\toprule
Model & Training & $r$ & $\Delta_\text{pert}$@55\% & Verdict \\
\midrule
Llama-3.1-8B-Instruct & Instruction & 0.471 & +6.6\% & PASS \\
DeepSeek-R1-Distill-Llama-8B & Reasoning & 0.19 & $-$0.8\% & FAIL \\
\midrule
Qwen2-7B-Instruct & Instruction & 0.22 & +1.2\% & FAIL \\
DeepSeek-R1-Distill-Qwen-7B & Reasoning & 0.147 & $-$2.2\% & FAIL \\
\bottomrule
\end{tabular}
\end{table}

DeepSeek-R1-Distill-Llama-8B shares Llama's GQA architecture but uses reasoning distillation training. The correlation drops from $r{=}0.47$ to $r{=}0.19$---a gap of 0.28 on the same architecture. Reasoning distillation introduces abundant meta-cognitive tokens (``wait'', ``let me reconsider'', ``alternatively'') whose KV representations do not follow the redundancy patterns of instruction-tuned content tokens. Critically, $\Delta_\text{pert}$ turns \emph{negative} ($-0.8\%$), indicating that high-similarity tokens in reasoning models are often meta-cognitive transition tokens whose MLP activations are essential for maintaining reasoning chains.

For the Qwen2 pair, both instruction and reasoning variants fail, but reasoning distillation further degrades the signal ($r$: $0.22 \rightarrow 0.147$, $\Delta_\text{pert}$: $+1.2\% \rightarrow -2.2\%$).

The implication is stark: the positive result on Llama-3.1-8B-Instruct reflects that specific model's training recipe, not a generic property of GQA attention.

\subsection{Sparsity Level Sweep on Llama}

Even on the one model where the signal exists, its practical utility is limited.

\begin{table}[h]
\centering
\caption{Cross-axis signal across sparsity levels on Llama-3.1-8B-Instruct.}
\label{tab:sparsity}
\small
\begin{tabular}{@{}ccccc@{}}
\toprule
Sparsity & $r$ & $\Delta_\text{pert}$ (\%) & Acc (uniform) & Acc (guided) \\
\midrule
50\% & 0.46 & +5.8 & 0.52 & 0.55 \\
55\% & 0.47 & +6.5 & 0.46 & 0.49 \\
60\% & 0.43 & +5.2 & 0.38 & 0.40 \\
65\% & 0.38 & +4.1 & 0.28 & 0.29 \\
70\% & 0.29 & +1.8 & 0.18 & 0.20 \\
\bottomrule
\end{tabular}
\end{table}

$\Delta_\text{pert}$ peaks at 55\% (+6.5\%) and decays monotonically. Beyond 65\%, the model collapses: both uniform and guided accuracy fall below 0.30 (baseline: 0.65), and the guided advantage shrinks to 2 percentage points at 70\%. The correlation $r$ itself degrades from 0.47 to 0.29, as extreme sparsity introduces noise that overwhelms the signal. The practical operating range is restricted to 50--60\%, a narrow window that limits deployment utility.

\subsection{Per-Layer Signal Decomposition}

\begin{table}[h]
\centering
\caption{Per-layer signal distribution on Llama-3.1-8B-Instruct (32 layers).}
\label{tab:perlayer}
\small
\begin{tabular}{@{}lcccc@{}}
\toprule
Layer Group & Layers & Mean $r$ & Mean $\Delta_\text{pert}$@55\% & Signal \\
\midrule
Early & 0--7 & 0.15 & +1.2\% & Weak/None \\
Early-Mid & 8--11 & 0.32 & +3.8\% & Moderate \\
Middle & 12--18 & 0.62 & +11.4\% & Strong \\
Late-Mid & 19--25 & 0.38 & +5.1\% & Moderate \\
Late & 26--31 & 0.18 & +2.0\% & Weak \\
\bottomrule
\end{tabular}
\end{table}

The aggregate $r{=}0.47$ is dominated by middle layers 12--18 (mean $r{=}0.62$). Only 9 of 32 layers have $r{>}0.4$; 12 layers have $r{<}0.2$. The layer-level standard deviation of $r$ is 0.18 (median $r{=}0.31$), indicating high heterogeneity. If cross-axis coordination were an architectural property of GQA, we would expect a more uniform distribution. The concentration in middle layers suggests it is a byproduct of Llama's specific compositional-reasoning computation, where early layers handle token mixing (KV similarity does not reflect semantic redundancy) and late layers handle output projection (MLP activations are critical for all tokens).

\subsection{Per-Head Analysis}

\begin{table}[h]
\centering
\caption{Per-head KV contribution to cross-axis signal.}
\label{tab:perhead}
\small
\begin{tabular}{@{}lcccccccc@{}}
\toprule
& H0 & H1 & H2 & H3 & H4 & H5 & H6 & H7 \\
\midrule
Llama $r$ & 0.52 & 0.48 & 0.45 & 0.42 & 0.41 & 0.35 & 0.28 & 0.14 \\
\bottomrule
\end{tabular}
\vspace{2pt}
\begin{tabular}{@{}lcccc@{}}
\toprule
& H0 & H1 & H2 & H3 \\
\midrule
DeepSeek $r$ & 0.15 & 0.12 & 0.09 & 0.08 \\
\bottomrule
\end{tabular}
\end{table}

Within Llama, per-head signal varies substantially: $r$ ranges from 0.52 (H0) to 0.14 (H7), with mean 0.41 and std 0.12. Five of eight heads show $r{>}0.4$, but one head (H7) has essentially no signal. This intra-model heterogeneity further supports the conclusion that the signal arises from specific learned patterns rather than an inherent architectural property. DeepSeek's four KV heads all show $r{\leq}0.15$ (mean 0.11, std 0.06).

\subsection{Downstream Two-Tier Sparsity}

\begin{table}[h]
\centering
\caption{Two-tier guided sparsity vs.\ uniform sparsity on Llama-3.1-8B-Instruct.}
\label{tab:twotier}
\small
\begin{tabular}{@{}lccc@{}}
\toprule
Strategy & Avg Sparsity & Acc & PPL \\
\midrule
No sparsity (baseline) & 0\% & 0.65 & 6.8 \\
Uniform 48\% & 48\% & 0.49 & 7.8 \\
Guided 55\%/40\% & $\sim$48\% & 0.50 & 7.7 \\
$\Delta$ (guided $-$ uniform) & --- & +0.01 & $-$0.1 \\
\bottomrule
\end{tabular}
\end{table}

The downstream test is the most telling result. At matched average sparsity ($\sim$48\%), guided two-tier allocation (55\% on high-similarity tokens, 40\% on low-similarity tokens) yields $\Delta_\text{acc}{=}{+}0.01$---statistically insignificant given a binomial 95\% CI of ${\pm}0.07$ on 20 problems. Despite $r{=}0.47$ in the correlation analysis, the cross-axis signal does not translate to meaningful practical benefit. The per-problem variance and layer concentration of the signal are absorbed into noise at the task level.

\subsection{CARC-KV: O(1) KV Eviction}

Independent of the cross-axis narrative, the EMA anchor similarity mechanism provides a viable $O(1)$ KV eviction method.

\begin{table}[h]
\centering
\caption{KV eviction accuracy by budget on Llama-3.1-8B-Instruct.}
\label{tab:kveviction}
\small
\begin{tabular}{@{}lcccc@{}}
\toprule
Method & Complexity & 30\% Budget & 50\% Budget & 70\% Budget \\
\midrule
Full KV & --- & --- & --- & 0.65 \\
CARC-KV & $O(1)$ & 0.42 & 0.58 & 0.63 \\
H2O & $O(n)$ & 0.44 & 0.59 & 0.64 \\
SnapKV & $O(n)$ & 0.46 & 0.60 & 0.64 \\
StreamingLLM & $O(1)$ & 0.31 & 0.45 & 0.55 \\
\bottomrule
\end{tabular}
\end{table}

CARC-KV trails H2O by 1--2 percentage points and SnapKV by 2--4 percentage points in accuracy, but significantly outperforms StreamingLLM (the other $O(1)$ method) by 8--13 percentage points. The key advantage is computational: CARC-KV's per-step eviction cost is constant regardless of sequence length.

\begin{table}[h]
\centering
\caption{Per-step eviction overhead (milliseconds).}
\label{tab:wallclock}
\small
\begin{tabular}{@{}cccc@{}}
\toprule
Seq Length & CARC-KV (ms) & H2O (ms) & Speedup \\
\midrule
4K & 0.8 & 2.1 & 2.6$\times$ \\
8K & 0.9 & 3.8 & 4.2$\times$ \\
16K & 1.0 & 7.2 & 7.2$\times$ \\
32K & 1.1 & 14.1 & 12.8$\times$ \\
\bottomrule
\end{tabular}
\end{table}

At 32K tokens, CARC-KV achieves 12.8$\times$ speedup in eviction overhead versus H2O. CARC-KV's overhead remains nearly constant (0.8--1.1ms) while H2O scales linearly (2.1--14.1ms). At 50\% KV budget, the accuracy--perplexity tradeoff is: CARC-KV achieves accuracy 0.58 / PPL 8.2 versus H2O's 0.59 / PPL 7.9---a modest gap in exchange for an order-of-magnitude reduction in eviction cost.

\subsection{Ablation: Key-Only vs.\ Key+Value Anchor}

\begin{table}[h]
\centering
\caption{Anchor vector ablation at 50\% KV budget.}
\label{tab:ablation}
\small
\begin{tabular}{@{}lcc@{}}
\toprule
Anchor Type & Acc@50\% & PPL@50\% \\
\midrule
Key-only & 0.56 & 8.6 \\
Key+Value & 0.58 & 8.2 \\
$\Delta$ & +0.02 & $-$0.4 \\
\bottomrule
\end{tabular}
\end{table}

Using concatenated key+value vectors for the EMA anchor provides a marginal improvement ($\Delta_\text{acc}{=}{+}0.02$, $\Delta_\text{PPL}{=}{-}0.4$). The small gap indicates the method is robust to this design choice, as key vectors already capture most token-identity information in GQA.

\section{Discussion}

\subsection{Why Does Only Llama Show the Signal?}

Our results rule out architecture (GQA) as the explanation, since Mistral-7B (same GQA configuration) and Qwen2-7B (GQA with different head count) both fail. The confound control further rules out architecture by showing that reasoning distillation destroys the signal on the same Llama architecture.

The most parsimonious explanation is that Llama-3.1-8B-Instruct's specific combination of pretraining data distribution, instruction tuning recipe, and full attention produces a token distribution where content redundancy happens to align with MLP sparsity tolerance. This alignment is an idiosyncrasy, not a generalizable principle.

Contributing factors that may explain cross-model differences include:
\begin{itemize}
    \item \textbf{Attention mechanism}: Mistral's SWA limits attention span to 4096 tokens, causing the EMA anchor to lose representativeness for late-sequence tokens.
    \item \textbf{KV head count}: Qwen2's 4 KV heads (vs.\ Llama's 8) force more generic KV representations with higher group ratio.
    \item \textbf{Training paradigm}: Reasoning distillation introduces meta-cognitive tokens whose high KV similarity does not imply low information content.
\end{itemize}

However, these factors interact in complex ways that resist clean isolation.

\subsection{Methodological Contribution}

While the cross-axis hypothesis fails as a general principle, the testing protocol itself has value:
\begin{enumerate}
    \item \textbf{Multi-model replication}: Testing on $\geq$3 models with controlled confounds before claiming architectural properties.
    \item \textbf{Per-layer/per-head decomposition}: Revealing whether aggregate statistics mask heterogeneous per-component behavior.
    \item \textbf{Downstream validation}: Requiring practical task-level improvement, not just correlation significance.
    \item \textbf{Confound control}: Varying training paradigm on the same architecture to separate architecture from training effects.
\end{enumerate}

We recommend this protocol for future work on cross-component efficiency signals.

\subsection{CARC-KV as Standalone Contribution}

The $O(1)$ KV eviction mechanism stands independent of the cross-axis narrative. CARC-KV's EMA anchor similarity requires one dot product per step regardless of cache size, achieving accuracy within 1--2\% of $O(n)$ methods while providing 4--13$\times$ speedup at practical sequence lengths (8K--32K). For latency-sensitive deployment scenarios where eviction overhead is a bottleneck, this tradeoff is favorable.

\section{Related Work}

\paragraph{KV cache eviction.} H2O \citep{zhang2023h2o} maintains cumulative attention scores for importance-based eviction. SnapKV \citep{li2024snapkv} uses observation windows for attention pattern clustering. StreamingLLM \citep{xiao2024streamingllm} combines attention sinks with sliding windows. All prior $O(n)$ methods achieve higher accuracy but with per-step cost that scales with cache size. CARC-KV bridges the gap between StreamingLLM's $O(1)$ simplicity and attention-score methods' accuracy.

\paragraph{Contextual MLP sparsity.} Deja Vu \citep{liu2023dejavu} demonstrates that MLP activation sparsity is input-dependent and predictable. CATS \citep{lee2024cats} and PowerInfer \citep{song2024powerinfer} exploit this for inference acceleration. Our work attempted to unify MLP sparsity prediction with KV redundancy signals, finding this unification does not generalize.

\paragraph{Multi-axis efficiency.} Token Filtering \citep{kim2024tokenfiltering} and CLADA \citep{wang2024clada} coordinate multiple efficiency axes. Our negative result suggests caution when assuming signals transfer across axes without per-model validation.

\section{Conclusion}

We set out to test whether KV-cache redundancy signals can predict MLP sparsity tolerance across models---the cross-axis coordination hypothesis. Our five-model study produces a clear negative answer: the signal is model-specific, not architecture-generic. Only one of five tested models (Llama-3.1-8B-Instruct) exhibits meaningful cross-axis correlation, and even there, the downstream accuracy benefit is negligible ($+$0.01). Per-layer analysis reveals the signal concentrates in 7 of 32 layers, and confound controls show training paradigm matters more than attention architecture.

We report this negative result because it provides useful information for the community: researchers pursuing multi-axis efficiency coordination should validate signals across models before claiming generality, and should demand downstream task-level improvements rather than relying on correlation metrics alone.

As a positive contribution, CARC-KV demonstrates that $O(1)$ EMA-based KV eviction is competitive with $O(n)$ methods (within 1--2\% accuracy) while achieving up to 12.8$\times$ speedup in eviction overhead, offering a practical efficiency primitive independent of cross-axis coordination claims.

\subsection*{Limitations}

All models tested are in the 7--8B parameter range; larger models may exhibit different redundancy patterns. The evaluation uses 20 mathematical reasoning problems; broader task coverage would strengthen conclusions. CARC-KV wall-clock measurements reflect eviction overhead only, not total inference time. The EMA decay rate ($\alpha{=}0.1$) was not extensively tuned per model.

\bibliographystyle{plainnat}
\begin{thebibliography}{10}

\bibitem[Zhang et~al.(2023)]{zhang2023h2o}
Zhenyu Zhang, Ying Sheng, Tianyi Zhou, Tianlong Chen, Lianmin Zheng, Ruisi Cai, Zhao Song, Yuandong Tian, Christopher R{\'e}, Clark Barrett, Zhangyang Wang, and Beidi Chen.
\newblock {H2O}: Heavy-hitter oracle for efficient generative inference of large language models.
\newblock In \emph{NeurIPS}, 2023.

\bibitem[Li et~al.(2024)]{li2024snapkv}
Yuhong Li, Yingbing Huang, Bowen Yang, Bharat Venkitesh, Acyr Locatelli, Hanchen Ye, Tianle Cai, Patrick Lewis, and Deming Chen.
\newblock {SnapKV}: {LLM} knows what you are looking for before generation.
\newblock In \emph{NeurIPS}, 2024.

\bibitem[Liu et~al.(2023)]{liu2023dejavu}
Zichang Liu, Jue Wang, Tri Dao, Tianyi Zhou, Binhang Yuan, Zhao Song, Anshumali Shrivastava, Ce Zhang, Yuandong Tian, Christopher R{\'e}, and Beidi Chen.
\newblock Deja {V}u: Contextual sparsity for efficient {LLM}s at inference time.
\newblock In \emph{ICML}, 2023.

\bibitem[Xiao et~al.(2024)]{xiao2024streamingllm}
Guangxuan Xiao, Yuandong Tian, Beidi Chen, Song Han, and Mike Lewis.
\newblock Efficient streaming language models with attention sinks.
\newblock In \emph{ICLR}, 2024.

\bibitem[Lee et~al.(2024)]{lee2024cats}
Dongjin Lee, Xinyu Zhao, Fatemeh Asgarinejad, Nayoung Lee, Bingbing Li, and Changho Hwang.
\newblock {CATS}: Contextually-aware thresholding for sparsity in large language models.
\newblock \emph{arXiv preprint arXiv:2404.08763}, 2024.

\bibitem[Song et~al.(2024)]{song2024powerinfer}
Yixin Song, Zeyu Mi, Haotong Xie, and Haibo Chen.
\newblock {PowerInfer}: Fast large language model serving with a consumer-grade {GPU}.
\newblock In \emph{MLSys}, 2024.

\bibitem[Kim et~al.(2024)]{kim2024tokenfiltering}
Sehyun Kim et~al.
\newblock Token filtering for efficient inference.
\newblock \emph{arXiv preprint}, 2024.

\bibitem[Wang et~al.(2024)]{wang2024clada}
Xu Wang et~al.
\newblock {CLADA}: Contrastive learning aided data augmentation for efficient inference.
\newblock \emph{arXiv preprint}, 2024.

\end{thebibliography}

\end{document}
